\section{Discussion and Limitations}
\label{sec:discussion}
\paragraph{Claim scope.}
The fixed-byte studies establish useful regimes: better CIDEr than VL-Cache
(R), stronger tight-budget quality than LOOK-M (R), and lower HF eager mean
request time. Native PrefixKV improves reference-output fidelity, but unequal
bytes and decode maintenance prevent attribution solely to selection.
FullKV-reference ROUGE-L measures output agreement, not necessarily answer
correctness. Near-FullKV PPL does not guarantee free-running answer quality.

\paragraph{Memory boundaries.}
Active-cache savings do not imply equal process-HBM savings. All paths start
from full prefill; source allocations can remain during planning. Peak here
is live decode-cache storage at transaction boundaries, not transient process
peak. Append-only growth can dominate long outputs; ongoing eviction can
lower final memory while altering outputs. A smaller-at-equal-quality claim
needs a separate paired test and a named competitor. E6's strict test failed.

\paragraph{Generality and attribution.}
Main results use batch-one HF eager and two LLaVA checkpoints. Eager attention
materializes scores and softmax across positions; fused/paged kernels change
that cost \citep{dao2022flash,dao2023flash2,kwon2023paged}. We make no
serving-scale speed claim. Baseline ports and common-byte adaptations differ
from official native implementations; the four-example PrefixKV parity gate
is limited evidence. Its shared-byte collapse does not describe the native
paper setting.

\paragraph{Remaining research.}
MileBench covers four subtasks, with 96 or 128 applicable examples per metric.
Uncertainty remains substantial. Prior development use is disclosed; E6
outcomes do not retune the candidate. Future work should use larger untouched
cohorts and isolate routing, forecast, and value/coverage effects.
Prototype merging and residual repair require separate adapter parity and
quality studies; their specifications appear as extensions in the appendix.
